
\documentclass[11pt,a4paper]{ctexart}

\usepackage{amsmath,amssymb,bm}
\usepackage{geometry}
\usepackage{booktabs}
\usepackage{array}
\usepackage{tikz}
\usepackage{pgfplots}
\usepackage{hyperref}
\usepackage{xcolor}
\usepackage{enumitem}

\usetikzlibrary{arrows.meta,positioning,calc}
\usepgfplotslibrary{groupplots}
\pgfplotsset{compat=1.18}

\geometry{left=2.1cm,right=2.1cm,top=2.2cm,bottom=2.2cm}
\setlength{\parindent}{2em}
\setlength{\parskip}{0.35em}
\setlist[itemize]{itemsep=0.2em,topsep=0.3em}
\setlist[enumerate]{itemsep=0.2em,topsep=0.3em}

\hypersetup{
    colorlinks=true,
    linkcolor=black,
    citecolor=black,
    urlcolor=blue
}

\newcommand{\KL}[2]{D_{\mathrm{KL}}\!\left(#1\,\middle\|\,#2\right)}
\newcommand{\E}{\mathbb{E}}
\newcommand{\piold}{\pi_{\mathrm{old}}}
\newcommand{\piref}{\pi_{\mathrm{ref}}}
\newcommand{\pitheta}{\pi_{\theta}}

\newenvironment{intuition}
{\begin{quote}\small\textbf{直觉：}}
{\end{quote}}

\title{\textbf{KL 散度：从数学推导到 LLM 强化学习}}
\author{}
\date{}

\begin{document}

\maketitle
\vspace{-2em}

\begin{center}
\small
目标：看懂 KL 散度的数学结构，并理解它在 PPO、DPO、GRPO、On Policy Distillation 中究竟控制了什么。
\end{center}

\tableofcontents

\section{先给结论：KL 在 LLM 训练里做什么}

KL 散度衡量两个概率分布有多不一样。对大语言模型而言，一个模型在每个 token 上都会输出一个概率分布，所以两个模型之间的差异天然可以用 KL 来衡量。

在 LLM 强化学习与对齐中，KL 最常见的作用可以概括为一句话：

\begin{intuition}
奖励负责告诉模型“往哪里走”，KL 负责限制“不要一次走得太远”。
\end{intuition}

例如，一个 SFT 模型已经会正常回答问题。强化学习发现某类答案 reward 很高，如果完全只追 reward，模型可能快速把概率集中到一些取巧模式上。加入对 reference model 的 KL 惩罚，相当于要求：

\[
\text{新模型既要拿更高 reward，也要尽量保留原模型已有的语言能力与行为分布。}
\]

最典型的目标写成

\begin{equation}
\max_{\pi}
\quad
\E_{y\sim\pi(\cdot|x)}[r(x,y)]
-
\beta \KL{\pi(\cdot|x)}{\piref(\cdot|x)}.
\label{eq:regularized_rl}
\end{equation}

其中 $\beta>0$ 控制 KL 约束强度。

\begin{itemize}
    \item $\beta$ 小：允许模型更大胆地偏离 reference policy。
    \item $\beta$ 大：模型更新更保守，更接近 reference policy。
\end{itemize}

后面的 PPO、DPO、GRPO、OPD，虽然形式不同，但都能看到“控制分布偏移”这个核心思想。

\section{KL 散度的定义}

\subsection{离散分布}

设离散随机变量 $X$ 的两个概率分布为

\[
P=\{p_1,p_2,\ldots,p_n\},
\qquad
Q=\{q_1,q_2,\ldots,q_n\},
\]

满足

\[
p_i\ge 0,\qquad q_i\ge 0,
\qquad
\sum_{i=1}^{n}p_i=1,
\qquad
\sum_{i=1}^{n}q_i=1.
\]

KL 散度定义为

\begin{equation}
\boxed{
\KL{P}{Q}
=
\sum_{i=1}^{n}
p_i\log\frac{p_i}{q_i}
}
\label{eq:kl_discrete}
\end{equation}

其中通常使用自然对数，因此单位是 nat。

也可以把求和写成期望：

\begin{align}
\KL{P}{Q}
&=
\sum_i p_i\log\frac{p_i}{q_i}
\\
&=
\E_{x\sim P}
\left[
\log\frac{P(x)}{Q(x)}
\right].
\label{eq:kl_expectation}
\end{align}

因此 KL 的本质就是：

\[
\boxed{
\text{在 }P\text{ 真正会产生的数据上，平均比较 }P\text{ 与 }Q\text{ 的 log probability 差异。}
}
\]

\subsection{连续分布}

若 $P,Q$ 有概率密度 $p(x),q(x)$，则

\begin{equation}
\KL{P}{Q}
=
\int
p(x)\log\frac{p(x)}{q(x)}\,dx.
\end{equation}

离散情形中的求和被积分替代，数学含义保持一致。

\subsection{为什么是 log probability ratio}

考虑某个事件 $x$：

\[
\log\frac{P(x)}{Q(x)}
=
\log P(x)-\log Q(x).
\]

若 $P(x)>Q(x)$，则

\[
\log\frac{P(x)}{Q(x)}>0.
\]

说明 $Q$ 低估了这个在 $P$ 中较常见的事件。

若 $P(x)<Q(x)$，则对应项为负。

KL 再用 $P(x)$ 作为权重求平均：

\[
\sum_x P(x)\log\frac{P(x)}{Q(x)}.
\]

所以，$P$ 经常出现的位置会得到更大的权重。

\begin{intuition}
把 $P$ 想成“真实分布”，把 $Q$ 想成“模型分布”。KL 会重点检查真实世界经常出现的情况，模型有没有给出足够概率。
\end{intuition}

\section{从交叉熵一步一步推导 KL}

定义熵

\begin{equation}
H(P)
=
-\sum_i p_i\log p_i.
\end{equation}

定义交叉熵

\begin{equation}
H(P,Q)
=
-\sum_i p_i\log q_i.
\end{equation}

计算两者之差：

\begin{align}
H(P,Q)-H(P)
&=
\left(-\sum_i p_i\log q_i\right)
-
\left(-\sum_i p_i\log p_i\right)
\\
&=
-\sum_i p_i\log q_i
+
\sum_i p_i\log p_i
\\
&=
\sum_i p_i
\left(
\log p_i-\log q_i
\right)
\\
&=
\sum_i p_i
\log\frac{p_i}{q_i}
\\
&=
\KL{P}{Q}.
\end{align}

因此

\begin{equation}
\boxed{
H(P,Q)=H(P)+\KL{P}{Q}
}
\end{equation}

或者

\begin{equation}
\boxed{
\KL{P}{Q}=H(P,Q)-H(P).
}
\end{equation}

如果 $P$ 固定，则 $H(P)$ 是常数，所以

\[
\min_Q H(P,Q)
\quad\Longleftrightarrow\quad
\min_Q \KL{P}{Q}.
\]

这就是最大似然训练和 forward KL 之间的重要联系。

\begin{intuition}
交叉熵等于“数据本身 unavoidable 的不确定性”加上“模型分布不准确额外付出的代价”。后面这部分就是 KL。
\end{intuition}

\section{KL 为什么一定非负}

KL 满足

\begin{equation}
\boxed{
\KL{P}{Q}\ge 0.
}
\end{equation}

下面不跳步证明。

假设 $p_i>0,q_i>0$。由 Jensen 不等式，因为 $-\log x$ 是凸函数，

\begin{equation}
\E[-\log X]
\ge
-\log\E[X].
\end{equation}

令

\[
X=\frac{q_i}{p_i},
\]

并让下标 $i$ 按照分布 $P$ 采样，则

\begin{align}
\KL{P}{Q}
&=
\sum_i p_i\log\frac{p_i}{q_i}
\\
&=
-\sum_i p_i\log\frac{q_i}{p_i}
\\
&=
\E_{i\sim P}
\left[
-\log\frac{q_i}{p_i}
\right]
\\
&\ge
-\log
\E_{i\sim P}
\left[
\frac{q_i}{p_i}
\right].
\end{align}

继续计算期望：

\begin{align}
\E_{i\sim P}
\left[
\frac{q_i}{p_i}
\right]
&=
\sum_i
p_i\frac{q_i}{p_i}
\\
&=
\sum_i q_i
\\
&=
1.
\end{align}

因此

\begin{align}
\KL{P}{Q}
&\ge
-\log 1
\\
&=
0.
\end{align}

所以

\[
\boxed{
\KL{P}{Q}\ge0.
}
\]

Jensen 取等号要求

\[
\frac{q_i}{p_i}=c
\]

对所有 $i$ 为常数。又因为

\[
\sum_i p_i=\sum_i q_i=1,
\]

所以 $c=1$，于是

\[
p_i=q_i,\quad \forall i.
\]

故

\begin{equation}
\boxed{
\KL{P}{Q}=0
\iff
P=Q.
}
\end{equation}

\section{KL 为什么不对称}

一般有

\[
\KL{P}{Q}
\neq
\KL{Q}{P}.
\]

例如

\[
P=(0.9,0.1),
\qquad
Q=(0.5,0.5).
\]

先算 $\KL{P}{Q}$：

\begin{align}
\KL{P}{Q}
&=
0.9\log\frac{0.9}{0.5}
+
0.1\log\frac{0.1}{0.5}
\\
&=
0.9\log 1.8
+
0.1\log 0.2
\\
&\approx
0.368.
\end{align}

反过来：

\begin{align}
\KL{Q}{P}
&=
0.5\log\frac{0.5}{0.9}
+
0.5\log\frac{0.5}{0.1}
\\
&=
0.5\log\frac{5}{9}
+
0.5\log 5
\\
&\approx
0.511.
\end{align}

两者明显不同。

原因来自权重：

\[
\KL{P}{Q}
=
\E_{x\sim P}
\left[
\log\frac{P(x)}{Q(x)}
\right]
\]

主要关心 $P$ 经常访问的位置，而

\[
\KL{Q}{P}
=
\E_{x\sim Q}
\left[
\log\frac{Q(x)}{P(x)}
\right]
\]

主要关心 $Q$ 经常访问的位置。

\section{Forward KL 与 Reverse KL}

假设目标分布是 $P$，可学习模型是 $Q_\theta$。

通常称

\begin{equation}
\text{Forward KL}
=
\KL{P}{Q_\theta},
\end{equation}

\begin{equation}
\text{Reverse KL}
=
\KL{Q_\theta}{P}.
\end{equation}

它们的行为差异非常重要。

\subsection{Forward KL 的倾向}

Forward KL 为

\[
\KL{P}{Q_\theta}
=
\E_{x\sim P}
\left[
\log\frac{P(x)}{Q_\theta(x)}
\right].
\]

如果某处

\[
P(x)>0,\qquad Q_\theta(x)\approx0,
\]

则

\[
\log\frac{P(x)}{Q_\theta(x)}
\rightarrow +\infty.
\]

所以 forward KL 很怕模型漏掉 $P$ 的高概率区域。

因此常见直觉是：

\[
\boxed{
\text{Forward KL 更偏向 mode covering。}
}
\]

\subsection{Reverse KL 的倾向}

Reverse KL 为

\[
\KL{Q_\theta}{P}
=
\E_{x\sim Q_\theta}
\left[
\log\frac{Q_\theta(x)}{P(x)}
\right].
\]

它主要惩罚模型把概率放到目标分布 $P$ 很低的位置。

如果 $P$ 有多个模式，而 $Q_\theta$ 的容量有限，$Q_\theta$ 可能更愿意集中到其中一个高概率模式。

因此常见直觉是：

\[
\boxed{
\text{Reverse KL 更偏向 mode seeking。}
}
\]

\begin{figure}[htbp]
\centering
\begin{tikzpicture}
\begin{groupplot}[
    group style={group size=3 by 1,horizontal sep=0.65cm},
    width=0.31\textwidth,
    height=0.23\textwidth,
    domain=-6:6,
    samples=180,
    axis lines=left,
    ticks=none,
    ymin=0,
    ymax=0.45,
    xlabel={$x$}
]
\nextgroupplot[title={目标分布 $P$}]
\addplot[thick]
{0.5*exp(-((x+2)^2)/2)/sqrt(2*pi)+0.5*exp(-((x-2)^2)/2)/sqrt(2*pi)};

\nextgroupplot[title={Forward KL 直觉}]
\addplot[thick,dashed]
{0.5*exp(-((x+2)^2)/2)/sqrt(2*pi)+0.5*exp(-((x-2)^2)/2)/sqrt(2*pi)};
\addplot[thick]
{exp(-(x^2)/(2*2.2^2))/(2.2*sqrt(2*pi))};

\nextgroupplot[title={Reverse KL 直觉}]
\addplot[thick,dashed]
{0.5*exp(-((x+2)^2)/2)/sqrt(2*pi)+0.5*exp(-((x-2)^2)/2)/sqrt(2*pi)};
\addplot[thick]
{exp(-((x-2)^2)/2)/sqrt(2*pi)};
\end{groupplot}
\end{tikzpicture}
\caption{模型容量受限时的典型直觉。虚线为目标分布，实线为近似分布。Forward KL 倾向覆盖多个模式，Reverse KL 倾向集中在一个高概率模式。}
\end{figure}

这里需要注意，“mode covering”和“mode seeking”是典型行为描述，并非对所有参数化分布都必然成立。

\section{从单个 token 推到整个 LLM 序列}

这是 KL 在 LLM 中最关键的一步。

给定 prompt $x$，模型生成序列

\[
y=(y_1,y_2,\ldots,y_T).
\]

自回归模型满足

\begin{equation}
\pitheta(y|x)
=
\prod_{t=1}^{T}
\pitheta(y_t|x,y_{<t}).
\end{equation}

reference model 同样满足

\begin{equation}
\piref(y|x)
=
\prod_{t=1}^{T}
\piref(y_t|x,y_{<t}).
\end{equation}

于是两者 sequence probability ratio 为

\begin{align}
\frac{\pitheta(y|x)}{\piref(y|x)}
&=
\frac{
\prod_{t=1}^{T}\pitheta(y_t|x,y_{<t})
}{
\prod_{t=1}^{T}\piref(y_t|x,y_{<t})
}
\\
&=
\prod_{t=1}^{T}
\frac{
\pitheta(y_t|x,y_{<t})
}{
\piref(y_t|x,y_{<t})
}.
\end{align}

取对数：

\begin{align}
\log
\frac{\pitheta(y|x)}{\piref(y|x)}
&=
\log
\prod_{t=1}^{T}
\frac{
\pitheta(y_t|x,y_{<t})
}{
\piref(y_t|x,y_{<t})
}
\\
&=
\sum_{t=1}^{T}
\log
\frac{
\pitheta(y_t|x,y_{<t})
}{
\piref(y_t|x,y_{<t})
}.
\label{eq:sequence_log_ratio}
\end{align}

再对 $y\sim\pitheta$ 取期望：

\begin{align}
\KL{\pitheta(\cdot|x)}{\piref(\cdot|x)}
&=
\E_{y\sim\pitheta}
\left[
\log
\frac{\pitheta(y|x)}{\piref(y|x)}
\right]
\\
&=
\E_{y\sim\pitheta}
\left[
\sum_{t=1}^{T}
\log
\frac{
\pitheta(y_t|x,y_{<t})
}{
\piref(y_t|x,y_{<t})
}
\right].
\label{eq:sequence_kl}
\end{align}

进一步使用条件 KL 的 chain rule：

\begin{align}
\KL{\pitheta(\cdot|x)}{\piref(\cdot|x)}
=
\sum_{t=1}^{T}
\E_{y_{<t}\sim\pitheta}
\left[
\KL{
\pitheta(\cdot|x,y_{<t})
}{
\piref(\cdot|x,y_{<t})
}
\right].
\label{eq:token_chain_rule}
\end{align}

这意味着：

\begin{equation}
\boxed{
\text{序列级 KL 可以理解为每一个生成位置上的 token 分布 KL 累积。}
}
\end{equation}

\begin{intuition}
一句回答有 500 个 token。模型每一步都可能比 reference model 偏一点。KL 会把这些局部偏移累积起来，因此长回答也需要特别注意 KL 的尺度与长度归一化。
\end{intuition}

\section{KL 的梯度：为什么它像一个“反向奖励”}

令 reference distribution $q(x)$ 固定，可训练分布为 $p_\theta(x)$：

\begin{equation}
D(\theta)
=
\sum_x
p_\theta(x)
\log
\frac{p_\theta(x)}{q(x)}.
\end{equation}

对 $\theta$ 求梯度：

\begin{align}
\nabla_\theta D(\theta)
&=
\sum_x
\nabla_\theta
\left[
p_\theta(x)
\log
\frac{p_\theta(x)}{q(x)}
\right]
\\
&=
\sum_x
\nabla_\theta p_\theta(x)
\log
\frac{p_\theta(x)}{q(x)}
+
\sum_x
p_\theta(x)
\nabla_\theta
\log p_\theta(x).
\end{align}

利用

\[
p_\theta(x)\nabla_\theta\log p_\theta(x)
=
\nabla_\theta p_\theta(x),
\]

得到

\begin{align}
\nabla_\theta D(\theta)
&=
\sum_x
\nabla_\theta p_\theta(x)
\log
\frac{p_\theta(x)}{q(x)}
+
\sum_x
\nabla_\theta p_\theta(x)
\\
&=
\sum_x
\nabla_\theta p_\theta(x)
\log
\frac{p_\theta(x)}{q(x)}
+
\nabla_\theta
\sum_x p_\theta(x).
\end{align}

因为

\[
\sum_x p_\theta(x)=1,
\]

所以

\[
\nabla_\theta
\sum_x p_\theta(x)
=
\nabla_\theta 1
=
0.
\]

因此

\begin{align}
\nabla_\theta D(\theta)
&=
\sum_x
\nabla_\theta p_\theta(x)
\log
\frac{p_\theta(x)}{q(x)}
\\
&=
\sum_x
p_\theta(x)
\nabla_\theta\log p_\theta(x)
\log
\frac{p_\theta(x)}{q(x)}
\\
&=
\E_{x\sim p_\theta}
\left[
\nabla_\theta\log p_\theta(x)
\log
\frac{p_\theta(x)}{q(x)}
\right].
\label{eq:kl_gradient}
\end{align}

考虑 KL 正则化强化学习目标

\[
J(\theta)
=
\E_{y\sim\pitheta}[r(y)]
-
\beta
\KL{\pitheta}{\piref}.
\]

忽略 baseline 等技巧，其 policy gradient 可以写成

\begin{equation}
\nabla_\theta J
=
\E_{y\sim\pitheta}
\left[
\nabla_\theta\log\pitheta(y)
\left(
r(y)
-
\beta
\log\frac{\pitheta(y)}{\piref(y)}
\right)
\right].
\end{equation}

因此可以把

\begin{equation}
-\beta
\log\frac{\pitheta(y)}{\piref(y)}
\end{equation}

看成一种额外的 reward correction。

\begin{intuition}
如果新模型把某条回答的概率抬得远高于 reference model，那么 log ratio 变大，KL correction 就会降低这条回答的有效奖励，从而阻止概率无限膨胀。
\end{intuition}

\section{KL 正则化 RL 的最优策略}

这一节非常重要，因为 DPO 的公式直接从这里来。

考虑固定 prompt $x$，省略 $x$：

\begin{equation}
\max_{\pi}
\left[
\sum_y \pi(y)r(y)
-
\beta
\sum_y
\pi(y)
\log
\frac{\pi(y)}{\piref(y)}
\right]
\end{equation}

约束为

\[
\sum_y\pi(y)=1.
\]

加入拉格朗日乘子 $\lambda$：

\begin{align}
\mathcal{L}(\pi,\lambda)
&=
\sum_y\pi(y)r(y)
-
\beta
\sum_y
\pi(y)
\log\frac{\pi(y)}{\piref(y)}
\\
&\quad
+
\lambda
\left(
\sum_y\pi(y)-1
\right).
\end{align}

对某个 $\pi(y)$ 求偏导：

\begin{align}
\frac{\partial\mathcal{L}}{\partial\pi(y)}
&=
r(y)
-
\beta
\frac{\partial}{\partial\pi(y)}
\left[
\pi(y)
\log\frac{\pi(y)}{\piref(y)}
\right]
+
\lambda.
\end{align}

使用

\[
\frac{d}{dz}
\left(
z\log z
\right)
=
\log z+1,
\]

有

\begin{align}
\frac{\partial}{\partial\pi(y)}
\left[
\pi(y)
\log\frac{\pi(y)}{\piref(y)}
\right]
&=
\log\frac{\pi(y)}{\piref(y)}
+1.
\end{align}

令偏导为零：

\begin{align}
0
&=
r(y)
-
\beta
\left(
\log\frac{\pi(y)}{\piref(y)}
+1
\right)
+
\lambda.
\end{align}

移项：

\begin{align}
\beta
\log\frac{\pi(y)}{\piref(y)}
&=
r(y)+\lambda-\beta.
\end{align}

除以 $\beta$：

\begin{align}
\log\frac{\pi(y)}{\piref(y)}
&=
\frac{r(y)}{\beta}
+
\frac{\lambda}{\beta}
-1.
\end{align}

两边取指数：

\begin{align}
\frac{\pi(y)}{\piref(y)}
&=
\exp
\left(
\frac{r(y)}{\beta}
\right)
\exp
\left(
\frac{\lambda}{\beta}-1
\right).
\end{align}

因此

\begin{equation}
\pi(y)
=
C
\piref(y)
\exp
\left(
\frac{r(y)}{\beta}
\right),
\end{equation}

其中 $C$ 与 $y$ 无关。

利用概率归一化：

\begin{align}
1
&=
\sum_y\pi(y)
\\
&=
C
\sum_y
\piref(y)
\exp
\left(
\frac{r(y)}{\beta}
\right).
\end{align}

定义 partition function

\begin{equation}
Z
=
\sum_y
\piref(y)
\exp
\left(
\frac{r(y)}{\beta}
\right),
\end{equation}

于是

\[
C=\frac{1}{Z}.
\]

最终得到

\begin{equation}
\boxed{
\pi^*(y)
=
\frac{1}{Z}
\piref(y)
\exp
\left(
\frac{r(y)}{\beta}
\right)
}
\label{eq:optimal_policy}
\end{equation}

这条式子非常直观：

\[
\boxed{
\text{新策略}
=
\text{旧策略先验}
\times
\text{reward 带来的指数加权}.
}
\]

\subsection{一个两答案的小例子}

设 reference policy 对两个回答 $A,B$ 的概率都是 $0.5$：

\[
\piref(A)=0.5,
\qquad
\piref(B)=0.5.
\]

奖励为

\[
r(A)=2,
\qquad
r(B)=0.
\]

由式 \eqref{eq:optimal_policy}：

\[
\frac{\pi^*(A)}{\pi^*(B)}
=
\frac{\piref(A)}{\piref(B)}
\exp
\left(
\frac{r(A)-r(B)}{\beta}
\right)
=
\exp\left(\frac{2}{\beta}\right).
\]

因此：

\begin{center}
\begin{tabular}{ccc}
\toprule
$\beta$ & $\pi^*(A)$ & 含义 \\
\midrule
$0.5$ & $\approx 0.982$ & reward 主导，更新激进 \\
$1$ & $\approx 0.881$ & 中等约束 \\
$2$ & $\approx 0.731$ & reference 约束更强 \\
\bottomrule
\end{tabular}
\end{center}

\section{PPO 中的 KL}

PPO 本身的核心是限制一次 policy update 的幅度。LLM RLHF 中又经常额外使用 reference model KL。因此需要区分两个对象：

\begin{enumerate}
    \item $\piold$：采样 rollout 时的旧策略，用 PPO ratio 与 clipping 控制单次更新。
    \item $\piref$：通常是冻结的 SFT 或初始策略，用 KL 防止长期漂移。
\end{enumerate}

\subsection{PPO ratio}

在 token $t$ 的状态 $s_t$ 和动作 $a_t$ 上定义

\begin{equation}
\rho_t(\theta)
=
\frac{
\pitheta(a_t|s_t)
}{
\piold(a_t|s_t)
}.
\end{equation}

PPO clipped objective 为

\begin{equation}
L_{\mathrm{clip}}
=
\E_t
\left[
\min
\left(
\rho_t A_t,
\operatorname{clip}
(\rho_t,1-\epsilon,1+\epsilon)A_t
\right)
\right].
\end{equation}

clipping 的作用是限制 $\pitheta$ 相对 $\piold$ 的单次变化。

\subsection{Reference KL}

RLHF 中常见的 reward 可以写成

\begin{equation}
r_{\mathrm{total}}
=
r_{\mathrm{RM}}
-
\beta
\log
\frac{
\pitheta(y_t|s_t)
}{
\piref(y_t|s_t)
}.
\end{equation}

沿整个序列累积：

\begin{equation}
R_{\mathrm{KL}}
=
-\beta
\sum_t
\log
\frac{
\pitheta(y_t|s_t)
}{
\piref(y_t|s_t)
}.
\end{equation}

在当前 policy 下取期望，就对应 sequence level forward KL。

\begin{figure}[htbp]
\centering
\begin{tikzpicture}[
    node distance=1.8cm,
    box/.style={draw,rounded corners,minimum width=3.1cm,minimum height=1.0cm,align=center},
    arr/.style={-{Latex[length=2.5mm]},thick}
]
\node[box] (ref) {Reference policy\\$\piref$};
\node[box,right=3.2cm of ref] (new) {Current policy\\$\pitheta$};
\node[box,above=1.4cm of new] (reward) {Reward model / verifier\\$r(x,y)$};

\draw[arr] (reward) -- node[right]{提高 reward} (new);
\draw[arr,<->] (ref) -- node[above]{KL 约束分布偏移} (new);
\end{tikzpicture}
\caption{LLM 强化学习中的典型结构。Reward 推动策略改变，KL 把策略拉回 reference policy 附近。}
\end{figure}

\begin{intuition}
PPO clipping 像“这一小步别跨太大”，reference KL 像“走很多步之后也别离出发点太远”。两种约束经常同时存在。
\end{intuition}

\section{DPO 中的 KL：损失里看不见，推导里一直存在}

DPO 省去了显式 reward model 训练与在线 PPO 循环，但它可以从 KL 正则化 RLHF 的最优策略直接推导出来。

由式 \eqref{eq:optimal_policy}：

\begin{equation}
\pi^*(y|x)
=
\frac{1}{Z(x)}
\piref(y|x)
\exp
\left(
\frac{r(x,y)}{\beta}
\right).
\end{equation}

移项：

\begin{align}
\frac{\pi^*(y|x)}{\piref(y|x)}
&=
\frac{1}{Z(x)}
\exp
\left(
\frac{r(x,y)}{\beta}
\right).
\end{align}

取对数：

\begin{align}
\log
\frac{\pi^*(y|x)}{\piref(y|x)}
&=
-\log Z(x)
+
\frac{r(x,y)}{\beta}.
\end{align}

乘以 $\beta$：

\begin{align}
\beta
\log
\frac{\pi^*(y|x)}{\piref(y|x)}
&=
-\beta\log Z(x)
+
r(x,y).
\end{align}

因此 reward 可以写成

\begin{equation}
\boxed{
r(x,y)
=
\beta
\log
\frac{\pi^*(y|x)}{\piref(y|x)}
+
\beta\log Z(x)
}
\label{eq:dpo_reward_parameterization}
\end{equation}

现在给定偏好数据

\[
(x,y_w,y_l),
\]

其中 $y_w$ 是 preferred response，$y_l$ 是 rejected response。

Bradley Terry 偏好模型写成

\begin{equation}
P(y_w\succ y_l|x)
=
\sigma
\left(
r(x,y_w)-r(x,y_l)
\right),
\end{equation}

其中

\[
\sigma(z)=\frac{1}{1+e^{-z}}.
\]

将式 \eqref{eq:dpo_reward_parameterization} 代入。

首先：

\begin{align}
r(x,y_w)
&=
\beta
\log
\frac{\pitheta(y_w|x)}{\piref(y_w|x)}
+
\beta\log Z(x),
\end{align}

\begin{align}
r(x,y_l)
&=
\beta
\log
\frac{\pitheta(y_l|x)}{\piref(y_l|x)}
+
\beta\log Z(x).
\end{align}

两式相减：

\begin{align}
r(x,y_w)-r(x,y_l)
&=
\beta
\log
\frac{\pitheta(y_w|x)}{\piref(y_w|x)}
\\
&\quad
-
\beta
\log
\frac{\pitheta(y_l|x)}{\piref(y_l|x)}
\\
&\quad
+
\beta\log Z(x)
-
\beta\log Z(x)
\\
&=
\beta
\left[
\log
\frac{\pitheta(y_w|x)}{\piref(y_w|x)}
-
\log
\frac{\pitheta(y_l|x)}{\piref(y_l|x)}
\right].
\end{align}

于是

\begin{equation}
P(y_w\succ y_l|x)
=
\sigma
\left(
\beta
\left[
\log
\frac{\pitheta(y_w|x)}{\piref(y_w|x)}
-
\log
\frac{\pitheta(y_l|x)}{\piref(y_l|x)}
\right]
\right).
\end{equation}

最大化偏好数据似然，就得到 DPO loss：

\begin{equation}
\boxed{
\mathcal{L}_{\mathrm{DPO}}
=
-
\E
\left[
\log
\sigma
\left(
\beta
\left[
\log
\frac{\pitheta(y_w|x)}{\piref(y_w|x)}
-
\log
\frac{\pitheta(y_l|x)}{\piref(y_l|x)}
\right]
\right)
\right].
}
\end{equation}

\begin{intuition}
DPO 在做的事情可以理解为：让模型相对 reference 更愿意产生 winner，同时相对 reference 更不愿意产生 loser。reference model 一直存在于概率比值中，所以 KL 正则化的思想已经被吸收到目标函数里。
\end{intuition}

DPO loss 中没有单独写

\[
-\beta\KL{\pitheta}{\piref},
\]

但 DPO 的闭式推导来自式 \eqref{eq:regularized_rl} 的 KL 正则化目标。因此说 DPO “完全与 KL 无关”是不准确的。

\section{GRPO 中的 KL}

GRPO 可以看成 PPO 风格的 policy optimization，但不依赖单独的 value model 来估计 advantage。对于同一个 prompt $q$，从旧策略采样一组回答：

\[
o_1,o_2,\ldots,o_G
\sim
\piold(\cdot|q).
\]

获得 rewards：

\[
r_1,r_2,\ldots,r_G.
\]

组内标准化得到 advantage：

\begin{equation}
A_i
=
\frac{
r_i-\operatorname{mean}(r_1,\ldots,r_G)
}{
\operatorname{std}(r_1,\ldots,r_G)
}.
\end{equation}

直觉上，同组里的回答彼此比较：

\begin{itemize}
    \item 比组平均好，$A_i>0$，提高其概率。
    \item 比组平均差，$A_i<0$，降低其概率。
\end{itemize}

token level ratio 为

\begin{equation}
\rho_{i,t}
=
\frac{
\pitheta(o_{i,t}|q,o_{i,<t})
}{
\piold(o_{i,t}|q,o_{i,<t})
}.
\end{equation}

一个典型 GRPO 目标写成

\begin{align}
J_{\mathrm{GRPO}}
=
\E
\Bigg[
\frac{1}{G}
\sum_{i=1}^{G}
\frac{1}{|o_i|}
\sum_t
\Big(
&
\min
[
\rho_{i,t}A_i,
\operatorname{clip}(\rho_{i,t},1-\epsilon,1+\epsilon)A_i
]
\nonumber\\
&
-
\beta D_{i,t}^{\mathrm{KL}}
\Big)
\Bigg].
\end{align}

DeepSeekMath 与 DeepSeek R1 相关公式中使用过如下单样本 KL estimator：

\begin{equation}
D_{i,t}^{\mathrm{KL}}
=
\frac{\piref(o_{i,t}|s_{i,t})}
{\pitheta(o_{i,t}|s_{i,t})}
-
\log
\frac{\piref(o_{i,t}|s_{i,t})}
{\pitheta(o_{i,t}|s_{i,t})}
-1.
\label{eq:grpo_kl_estimator}
\end{equation}

这个式子看起来和标准 KL 定义不同，但在 $o_{i,t}\sim\pitheta$ 时，它的期望正好对应 forward KL。

令

\[
R(a)
=
\frac{\piref(a|s)}{\pitheta(a|s)}.
\]

则

\begin{align}
\E_{a\sim\pitheta}
[
R(a)-\log R(a)-1
]
&=
\E_{\pitheta}[R(a)]
-
\E_{\pitheta}[\log R(a)]
-1.
\end{align}

第一项：

\begin{align}
\E_{\pitheta}[R(a)]
&=
\sum_a
\pitheta(a|s)
\frac{\piref(a|s)}{\pitheta(a|s)}
\\
&=
\sum_a\piref(a|s)
\\
&=
1.
\end{align}

第二项：

\begin{align}
-\E_{\pitheta}[\log R(a)]
&=
-\E_{\pitheta}
\left[
\log
\frac{\piref(a|s)}{\pitheta(a|s)}
\right]
\\
&=
\E_{\pitheta}
\left[
\log
\frac{\pitheta(a|s)}{\piref(a|s)}
\right]
\\
&=
\KL{\pitheta(\cdot|s)}{\piref(\cdot|s)}.
\end{align}

所以

\begin{align}
\E_{\pitheta}
[
R-\log R-1
]
&=
1
+
\KL{\pitheta}{\piref}
-
1
\\
&=
\KL{\pitheta}{\piref}.
\end{align}

因此式 \eqref{eq:grpo_kl_estimator} 仍然是在控制 current policy 与 reference policy 的偏离。

\begin{intuition}
GRPO 的主要变化在 advantage 怎么来。PPO 依赖 critic 估值，GRPO 让同一道题采样出来的一组回答互相比较。KL 的职责没有变，仍然负责防止 policy 为追求 reward 过度漂移。
\end{intuition}

\section{OPD：On Policy Distillation 中的 KL}

本文把 OPD 指 On Policy Distillation。

设 teacher 为

\[
p_T,
\]

student 为

\[
p_S^\theta.
\]

传统离线蒸馏通常在固定 teacher data 或固定训练数据上训练 student。问题是推理时 student 会生成自己的 token，一旦前面生成错，后续看到的 prefix 可能完全偏离训练集。

On Policy Distillation 的关键流程是：

\begin{enumerate}
    \item 给 student 一个 prompt $x$。
    \item student 自己生成轨迹 $y\sim p_S^\theta(\cdot|x)$。
    \item 在 student 真正访问到的每个 prefix $y_{<t}$ 上查询 teacher。
    \item 比较 teacher 与 student 的 next token distribution。
    \item 用 KL 或其他 divergence 更新 student。
\end{enumerate}

\begin{figure}[htbp]
\centering
\begin{tikzpicture}[
    node distance=1.35cm,
    box/.style={draw,rounded corners,minimum width=3.3cm,minimum height=0.9cm,align=center},
    arr/.style={-{Latex[length=2.5mm]},thick}
]
\node[box] (prompt) {Prompt $x$};
\node[box,right=1.6cm of prompt] (student) {Student 生成\\$y\sim p_S^\theta$};
\node[box,right=1.6cm of student] (prefix) {Student 实际访问的\\prefix $y_{<t}$};
\node[box,below=1.2cm of prefix] (teacher) {Teacher 给出\\next token distribution};
\node[box,left=1.6cm of teacher] (loss) {计算 divergence\\并更新 Student};

\draw[arr] (prompt) -- (student);
\draw[arr] (student) -- (prefix);
\draw[arr] (prefix) -- (teacher);
\draw[arr] (teacher) -- (loss);
\draw[arr] (loss) -- (student);
\end{tikzpicture}
\caption{On Policy Distillation 的基本闭环。Student 先暴露自己的真实错误状态，再由 Teacher 在这些状态上提供密集分布反馈。}
\end{figure}

\subsection{Forward KL 蒸馏}

按 teacher 为目标、student 为近似的常用记法：

\begin{equation}
D_{\mathrm{FKL}}
=
\KL{
p_T(\cdot|s_t)
}{
p_S^\theta(\cdot|s_t)
}.
\end{equation}

展开：

\begin{align}
D_{\mathrm{FKL}}
&=
\sum_v
p_T(v|s_t)
\log
\frac{
p_T(v|s_t)
}{
p_S^\theta(v|s_t)
}.
\end{align}

因为 teacher 固定，

\[
\sum_v p_T(v|s_t)\log p_T(v|s_t)
\]

对 student 参数是常数，所以最小化 forward KL 等价于最小化

\begin{equation}
-\sum_v
p_T(v|s_t)
\log p_S^\theta(v|s_t).
\end{equation}

这就是 soft target cross entropy。

\subsection{Reverse KL 蒸馏}

Reverse KL 为

\begin{equation}
D_{\mathrm{RKL}}
=
\KL{
p_S^\theta(\cdot|s_t)
}{
p_T(\cdot|s_t)
}.
\end{equation}

展开：

\begin{equation}
D_{\mathrm{RKL}}
=
\sum_v
p_S^\theta(v|s_t)
\log
\frac{
p_S^\theta(v|s_t)
}{
p_T(v|s_t)
}.
\end{equation}

它更强调：

\[
\text{Student 自己想输出的 token，Teacher 是否也认为合理。}
\]

如果 student 给某个 token 很高概率，而 teacher 给它很低概率，则

\[
\log
\frac{p_S^\theta(v|s_t)}{p_T(v|s_t)}
\]

会很大，从而产生强惩罚。

\begin{intuition}
Forward KL 更像“尽量把 teacher 会的各种可能性都学下来”。Reverse KL 更像“student 既然准备走这条路，就检查 teacher 是否认可这条路”。当 student 容量明显小于 teacher 时，后者有时更容易把有限容量集中到 teacher 的高概率模式。
\end{intuition}

GKD 工作把 on policy sampling 与多种 divergence 结合，并研究了 reverse KL 与 Jensen Shannon divergence 等选择。

\subsection{OPD 和强化学习的关系}

OPD 看起来像蒸馏，但和 RL 有一个很重要的共同点：

\[
\boxed{
\text{训练数据来自当前 student policy 自己访问到的状态。}
}
\]

这和 online RL 的 on policy 数据非常相似。

同时 teacher 给出的 token distribution 是稠密反馈：

\[
p_T(\cdot|s_t),
\]

它比一个只在整条回答结束后给出的 scalar reward

\[
r(x,y)\in\mathbb{R}
\]

信息更丰富。

因此可以粗略理解为：

\begin{itemize}
    \item SFT：给 student 一个示范答案。
    \item RL：告诉 student 整条轨迹最终得多少分。
    \item OPD：student 先自己走，然后 teacher 在 student 真正走到的每一步给出概率分布参考。
\end{itemize}

需要避免一个过强的说法：teacher 并没有逐 token 给出“绝对正确或错误”的标签。teacher 给的是一个概率分布，表示它在当前 prefix 下更倾向哪些 next token。

\section{四种方法放在一起比较}

\begin{center}
\small
\begin{tabular}{p{1.7cm}p{3.0cm}p{3.5cm}p{3.2cm}p{3.0cm}}
\toprule
方法 & 数据来源 & KL 或 reference 出现位置 & KL 主要作用 & 一句话理解 \\
\midrule
PPO
&
当前或旧 policy rollout
&
常显式约束 $\pitheta$ 与 $\piref$，同时 ratio clipping 约束 $\pitheta$ 与 $\piold$
&
限制 reward 驱动下的策略漂移
&
拿 reward，但别一步更新过猛，也别长期偏离 SFT 太远
\\
\addlinespace

DPO
&
离线 preference pairs
&
通过 $\log\pitheta-\log\piref$ 的偏好 log ratio 隐式继承 KL regularized RLHF
&
保持 reference 锚点，同时增大 winner 相对 loser 的偏好
&
直接学“更喜欢哪个回答”，省掉显式 PPO 循环
\\
\addlinespace

GRPO
&
当前或旧 policy 的 group rollout
&
显式 KL 到 reference，policy ratio 到 old policy
&
限制 group relative reward 导致的策略漂移
&
同一道题多采样几份，组内比高低，再保留 KL 安全绳
\\
\addlinespace

OPD
&
Student 自己生成的 on policy trajectories
&
Student 与 Teacher 的 token distribution divergence
&
在 student 真正访问的状态上纠正分布
&
先让 student 自己走，再让 teacher 在它走到的位置指导
\\
\bottomrule
\end{tabular}
\end{center}

\section{KL 系数 $\beta$ 到底控制什么}

回到

\[
J(\pi)
=
\E_\pi[r]
-
\beta
\KL{\pi}{\piref}.
\]

$\beta$ 是 reward 与“保持原模型行为”之间的交换系数。

\subsection{$\beta$ 太小}

如果

\[
\beta\rightarrow0,
\]

目标越来越接近

\[
\max_\pi\E_\pi[r].
\]

可能出现：

\begin{itemize}
    \item reward hacking；
    \item policy 快速塌到少量高 reward 模式；
    \item 语言质量或通用能力下降；
    \item 训练不稳定；
    \item exploration 过早消失。
\end{itemize}

\subsection{$\beta$ 太大}

如果 $\beta$ 很大，KL penalty 主导：

\[
\pi\approx\piref.
\]

可能出现：

\begin{itemize}
    \item 模型几乎学不到 reward 信号；
    \item reasoning 行为难以发生明显变化；
    \item policy improvement 过慢。
\end{itemize}

因此训练中经常关注：

\[
\text{reward},
\qquad
\text{KL},
\qquad
\text{entropy},
\qquad
\text{response length},
\]

而不是只看 reward。

\section{一个更深的理解：KL 是“更新预算”}

把 reference policy 想成模型已经掌握的行为分布。

reward 想把某些回答概率提高：

\[
\pitheta(y|x)\uparrow.
\]

但概率总和必须为 1：

\[
\sum_y\pitheta(y|x)=1.
\]

所以提高某些回答的概率，必然会压低另一些回答。

如果没有约束，模型为了快速提高 reward，可能重新分配大量概率质量。

KL 通过

\[
\log
\frac{
\pitheta(y|x)
}{
\piref(y|x)
}
\]

直接衡量这种概率质量重新分配。

因此可以把 KL 理解为一种 distribution shift budget：

\begin{equation}
\boxed{
\text{为了得到更高 reward，你愿意花多少“偏离原模型”的预算？}
}
\end{equation}

这也是为什么 KL 在 LLM 强化学习里如此常见。LLM 的 action space 是整个词表，trajectory 又很长，轻微的 token probability 改动经过多步自回归之后就可能产生很大的行为变化。

\section{常见误区}

\subsection{误区一：KL 是距离}

KL 满足非负性，但通常

\[
\KL{P}{Q}\neq\KL{Q}{P},
\]

因此它不是严格数学意义上的距离。

\subsection{误区二：PPO clipping 和 reference KL 是同一件事}

它们约束的对象通常不同：

\[
\pitheta
\leftrightarrow
\piold
\]

主要控制单次 policy update；

\[
\pitheta
\leftrightarrow
\piref
\]

主要控制长期相对初始模型的漂移。

\subsection{误区三：DPO 没写 KL，所以和 KL 没关系}

DPO loss 没有单独的 KL 项，但其核心 reward parameterization 来自 KL regularized RLHF 的闭式最优策略。

\subsection{误区四：GRPO 去掉 critic，所以也去掉 KL}

GRPO 的核心改动是使用 group relative reward 构造 advantage，从而省掉独立 value model。原始 DeepSeekMath 以及 DeepSeek R1 的相关 GRPO 形式仍包含 reference KL penalty。

\subsection{误区五：OPD 的 teacher 是逐 token 正误判定器}

更准确的说法是 teacher 提供 next token probability distribution。它给的是稠密偏好信号，而不是逐 token 的二元正确性标签。

\section{最终记忆框架}

如果只记住五句话：

\begin{enumerate}
    \item KL 衡量两个概率分布的差异：
    \[
    \KL{P}{Q}
    =
    \E_{P}
    \left[
    \log\frac{P}{Q}
    \right].
    \]

    \item KL 非负：
    \[
    \KL{P}{Q}\ge0,
    \qquad
    \KL{P}{Q}=0\iff P=Q.
    \]

    \item 对 LLM，自回归分解使 sequence KL 可以拆成每个 token 的 log probability ratio 累积。

    \item 在 RLHF 中：
    \[
    \text{reward 推动模型改变，KL 限制模型改变得太远。}
    \]

    \item PPO、DPO、GRPO、OPD 的实现路径不同，但 KL 都在处理同一个根问题：
    \[
    \boxed{
    \text{如何学习新的偏好，同时控制概率分布发生多大的变化。}
    }
    \]
\end{enumerate}

\section{参考文献}

\begin{thebibliography}{9}

\bibitem{kl1951}
S. Kullback and R. A. Leibler.
\newblock On Information and Sufficiency.
\newblock \emph{The Annals of Mathematical Statistics}, 1951.

\bibitem{ppo}
J. Schulman, F. Wolski, P. Dhariwal, A. Radford, and O. Klimov.
\newblock Proximal Policy Optimization Algorithms.
\newblock arXiv:1707.06347, 2017.
\newblock \url{https://arxiv.org/abs/1707.06347}

\bibitem{dpo}
R. Rafailov, A. Sharma, E. Mitchell, S. Ermon, C. D. Manning, and C. Finn.
\newblock Direct Preference Optimization: Your Language Model is Secretly a Reward Model.
\newblock arXiv:2305.18290, 2023.
\newblock \url{https://arxiv.org/abs/2305.18290}

\bibitem{deepseekmath}
Z. Shao et al.
\newblock DeepSeekMath: Pushing the Limits of Mathematical Reasoning in Open Language Models.
\newblock arXiv:2402.03300, 2024.
\newblock \url{https://arxiv.org/abs/2402.03300}

\bibitem{deepseekr1}
DeepSeek AI et al.
\newblock DeepSeek R1 Incentivizes Reasoning in LLMs through Reinforcement Learning.
\newblock \emph{Nature}, 2025.
\newblock \url{https://www.nature.com/articles/s41586-025-09422-z}

\bibitem{gkd}
R. Agarwal, N. Vieillard, Y. Zhou, P. Stanczyk, S. Ramos, M. Geist, and O. Bachem.
\newblock On Policy Distillation of Language Models: Learning from Self Generated Mistakes.
\newblock arXiv:2306.13649, 2023.
\newblock \url{https://arxiv.org/abs/2306.13649}

\end{thebibliography}

\end{document}
